\providecommand{\frG}{\mathfrak G}
\providecommand{\frL}{\mathfrak L}
\providecommand{\frM}{\mathfrak M}
\providecommand{\frN}{\mathfrak N}
\providecommand{\frT}{\mathfrak T}

\subsection*{\S 3. Modules relativement à un corps gauche}

Dans ce qui suit il ne s'agira que de représentations dans des corps gauches en général non commutatifs. On rassemblera donc quelques théorèmes simples sur les modules de formes linéaires sur un corps gauche.

\paragraph{1. Base normale d'un sous-module d'après une base donnée du module total.}
Soit \(A\) un corps gauche et soit
\[
        \frN=x_1A+\cdots+x_nA
\]
un module de formes linéaires de rang \(n\); soit en outre
\[
        \frL=z_1A+\cdots+z_\ell A
\]
un sous-module de rang \(\ell\le n\). Les \(z_i\) s'appellent base normale d'après \(x\) s'ils sont de la forme
\[
        z_i=x_i-(x_{\ell+1}\alpha_{i,\ell+1}+\cdots+x_n\alpha_{i,n}),
        \qquad i=1,\ldots,\ell,
\]
Tout module \(\frL\) possède, après numérotation convenable des \(x_i\), une base normale. En effet, après une numérotation convenable on obtient
\[
        \frN=\frL+x_{\ell+1}A+\cdots+x_nA,
\]
donc
\[
        x_i\equiv x_{\ell+1}\alpha_{i,\ell+1}+\cdots+x_n\alpha_{i,n}
        \pmod{\frL},\qquad i=1,\ldots,\ell. \tag{2}
\]
et par conséquent les
\[
        z_i=x_i-(x_{\ell+1}\alpha_{i,\ell+1}+\cdots+x_n\alpha_{i,n})
\]
correspondants appartiennent à \(\frL\) et épuisent \(\frL\), puisqu'ils forment, avec \(x_{\ell+1},\ldots,x_n\), une base de \(\frN\).

\paragraph{2. Module d'extension.}
Soit de nouveau \(A\) un corps gauche, et \(P\) un sous-corps gauche. Un \(A\)-module \(\frN\) de rang \(n\) s'appelle module d'extension d'un \(P\)-module \(\frM\); \(\frN=\frM_A\), si \(\frM\) est sous-module de même rang de \(\frN\). Si
\[
        \frM=y_1P+\cdots+y_nP,
\]
alors
\[
        \frN=\frM_A=y_1A+\cdots+y_nA.
\]
Réciproquement, à chaque \(P\)-module \(\frM=y_1P+\cdots+y_nP\) correspond, à isomorphie de modules près, un module d'extension \(\frM_A\) déterminé de façon univoque. Car le module \(\overline y_1A+\cdots+\overline y_nA\) contient un sous-module \(\overline{\frM}\) isomorphe à \(\frM\), qui peut être identifié à \(\frM\).

\paragraph{Conséquence a).}
Si \(z_1,\ldots,z_t\) sont des éléments de \(\frM\) linéairement indépendants relativement à \(P\), alors ils sont aussi linéairement indépendants sur \(A\) dans \(\frN=\frM_A\) (car ils peuvent être complétés en une base de \(\frM\), et donc de \(\frN\)). Tout \(P\)-module
\[
        \frT=z_1P+\cdots+z_tP
\]
contenu dans \(\frM\) engendre donc un module d'extension
\[
        \frT_A=z_1A+\cdots+z_tA\subset \frM_A.
\]

\paragraph{Conséquence b).}
Tout \(P\)-module \(\frT=z_1P+\cdots+z_tP\) contenu dans \(\frM\) est module de contraction, c'est-à-dire l'intersection de son module d'extension \(\frT_A\) avec \(\frM\):
\[
        \frT=\frT_A\cap\frM .
\]
Car \(\frT\subseteq \frT_A\cap\frM\); si, réciproquement, \(a\) appartient à \(\frT_A\cap\frM\), on a
\[
        a=\sum z_i\alpha_i,
\]
donc les éléments \(a,z_1,\ldots,z_t\) de \(\frM\) sont dépendants sur \(A\), donc d'après la conséquence a) aussi sur \(P\); par suite \(a\) appartient à \(\frT\).

\paragraph{3. Théorème sur les modules invariants.}
Soit de nouveau \(\frN=\frM_A\) module d'extension d'un \(P\)-module \(\frM\) de rang \(n\), et soit en outre \(P\) le corps gauche complet des invariants relativement à un groupe \(\frG\) d'automorphismes d'anneaux de \(A\). Le groupe \(\frG\) est défini comme domaine d'opérateurs de \(\frN=\frM_A\) par les prescriptions
\[
        Gm=m\quad\text{pour }m\text{ dans }\frM\text{ et }G\text{ dans }\frG
\]
et
\[
        G\sum m_i\alpha_i=\sum m_i\,G(\alpha_i)
        \quad\text{pour }m_i\text{ dans }\frM\text{ et }\alpha_i\text{ dans }A.
\]

\paragraph{Lemme.}
D'après ces prescriptions, \(\frM\) est constitué par la totalité des éléments de \(\frN\) qui restent fixes sous \(\frG\). En effet, si \(x_1,\ldots,x_n\) est une \(P\)-base de \(\frM\), donc \(w=x_1\alpha_1+\cdots+x_n\alpha_n\), alors \(G(w)=w\) pour tout \(G\) de \(\frG\) donne aussi \(G(\alpha_i)=\alpha_i\), et par conséquent, d'après l'hypothèse, \(\alpha_i\) appartient à \(P\).

\paragraph{Théorème.}
Les modules d'extension \(\frL_A\) des sous-modules \(\frL\) de \(\frM\), et eux seuls, sont des sous-modules admissibles relativement à \(\frG\) comme domaine d'opérateurs.

La première moitié est claire. Soit réciproquement \(\frL\) admissible relativement à \(\frG\), et soit \(z_1,\ldots,z_\ell\) une base normale de \(\frL\) (voir 1.) d'après \(x_1,\ldots,x_n\), où \(x_1,\ldots,x_n\) est choisi comme \(P\)-base de \(\frM\), donc est admis élément par élément par \(\frG\). Par hypothèse,
\[
        G(z_i)=x_i-\bigl(x_{\ell+1}G(\alpha_{i,\ell+1})
              +\cdots+x_nG(\alpha_{i,n})\bigr)
\]
est élément de \(\frL\) pour tout \(G\) de \(\frG\), donc, pour \(G\) fixe, exprimable linéairement par les \(z\); ainsi, comme il résulte de la comparaison des coefficients en \(x_1,\ldots,x_\ell\), \(G(z_i)=z_i\) pour tout \(G\) de \(\frG\). Les \(z\) appartiennent donc -- d'après le lemme -- à \(\frM\), et \(\frL\) devient extension du \(P\)-module
\[
        \frT=z_1P+\cdots+z_\ell P
\]
contenu dans \(\frM\), avec \(\frT=\frL\cap\frM\) (d'après 2).\footnote{Les théorèmes de ce paragraphe restent valables, au moyen de raisonnements simples de bon ordre, même lorsque le rang de \(\frM\) sur \(P\) devient infini (voir G. Kothe: Ein Beitrag zur Theorie der kommutativen Ringe ohne Endlichkeitsvoraussetzung \S 1, Gott. Nachr. 1931, p. 195--207).}

\clearpage
